\section{Related Work}\label{sec:related}

Our findings connect to several distinct research threads that have approached neural network dynamics from different angles.

\paragraph{Grokking and Algorithmic Learning.}
Power et al.~\cite{power2022grokking} discovered that small transformers trained on modular arithmetic exhibit delayed generalization—memorizing training data for thousands of steps before suddenly ``grokking'' the underlying algorithm. Nanda et al. provided mechanistic interpretations showing that grokking involves the emergence of Fourier-based circuits~\cite{elhage2021mathematical}. Our work reveals that this phase transition coincides with \textit{homeostatic crystallization}: not just algorithm discovery, but convergence to a precise geometric equilibrium. The 12-decimal precision we observe suggests grokking is better understood as a thermodynamic phase transition than as sudden circuit formation.

\paragraph{Information Bottleneck Theory.}
Tishby and Zaslavsky~\cite{tishby2015deep} proposed that deep learning involves an information bottleneck: networks compress input representations while preserving task-relevant information. Our EDI equilibrium represents a manifestation of this compression limit—the point where entropy reduction halts at a task-determined set point. Unlike prior IB work focusing on mutual information dynamics, we observe \textit{active maintenance} of this equilibrium via Le Chatelier compensation.

\paragraph{Superposition and Active Interference Management.}
Elhage et al.~\cite{elhage2022superposition} establish that neural networks represent more features than dimensions through ``superposition,'' a strategy requiring tolerance or filtering of ``interference'' between non-orthogonal feature representations. They identify discrete geometric configurations (polytopes) that models snap into via phase transitions, observe ``sticky'' dimensionality at specific fractions, and describe circuit motifs---including ``asymmetric superposition with inhibition''---designed to suppress interference and preserve feature fidelity. Our findings provide dynamical validation of this geometric picture. The 12-decimal equilibrium we observe (EDI = 0.611991..., corresponding to $\log(3)/\log(6)$) represents extreme precision convergence to exactly the kind of discrete configuration Elhage et al.\ describe. More critically, our Kill Test experiments (Section~\ref{sec:mechanism}) provide causal evidence that interference management is an \textit{active homeostatic process}, not merely a static architectural property. When we perturb one head (introducing interference), other heads immediately compensate to restore the equilibrium---the system actively defends its low-interference geometry. The layer-0 suppressors we identify in GPT-2 and Mistral (Section~\ref{sec:scale}) likely implement the same function as Elhage et al.'s ``inhibition'' neurons at production scale: filtering high-interference signals to preserve representational structure. Where Elhage et al.\ characterize the geometry of stable configurations, we characterize the dynamics by which systems reach and defend them.

\paragraph{Thermodynamic Framings of Learning.}
Kringelbach et al.~\cite{kringelbach2024thermodynamics} recently proposed ``thermodynamics of mind'' for biological neural systems, arguing that brains maintain far-from-equilibrium states through continuous energy expenditure. Friston's Free Energy Principle~\cite{friston2010free} frames perception as entropy minimization. We extend these insights to artificial systems: networks don't passively settle into loss minima—they actively regulate internal geometry via distributed compensation. The kill test (Section~\ref{sec:mechanism}) directly demonstrates this: perturb one component, and others adjust to preserve the equilibrium.

\paragraph{Singular Learning Theory.}
Watanabe's algebraic geometry framework~\cite{watanabe2009algebraic} characterizes loss landscapes through the Real Log Canonical Threshold (RLCT), predicting that models converge to singular points with specific geometric properties. Our forced attractor at EDI $\approx 0.44$–$0.46$ under homeostatic pressure may correspond to an RLCT-constrained basin—a point where the architecture's information capacity creates a geometric ceiling. The gap between the natural equilibrium (0.61) and the forced attractor (0.44) quantifies this constraint.

\paragraph{Dual-Timescale Learning.}
Complementary learning systems in neuroscience~\cite{mcclelland1995complementary} propose that fast (hippocampal) and slow (cortical) learning operate at different timescales to balance plasticity and stability. Our dual-timescale training (Section~\ref{sec:methods}) implements an analogous separation, but with an explicit homeostatic objective. Unlike biological models focused on memory consolidation, we show how timescale separation enables \textit{engineering control} over phase transition dynamics (up to 93\% acceleration of crystallization in preliminary experiments).

\paragraph{Attention Entropy and Head Specialization.}
Voita et al.~\cite{voita2019analyzing} showed that attention heads specialize during training, with many becoming ``dead'' or highly sparse. Our EDI metric captures this specialization, but reveals it is not random: heads converge to shared equilibria with extraordinary precision. This connects to work on attention entropy as a diagnostic~\cite{jain2019attention}, extending it from interpretation tool to fundamental conserved quantity.
